%Chargement des paquets
\usepackage{amsmath}
\usepackage{amsfonts}
\usepackage{amssymb}
\usepackage{enumerate}
\usepackage{mathtools}
\usepackage{bbm}
\usepackage{xparse, etoolbox}
\usepackage{enumerate}
\usepackage{enumitem}
\usepackage{mathabx}
\usepackage{babel}[french]

%environment
\usepackage{amsthm}
\usepackage{keytheorems}
\usepackage{hyperref} 

%Interface théorème
\newkeytheorem{lem}[
	name = Lemme
]

\newkeytheorem{prop}[
	name = Proposition
]

\newkeytheorem{thm}[
	name = Théorème
]

\newkeytheorem{coro}[
	name = Corollaire
]

\renewcommand*{\proofname}{Démonstration}

\theoremstyle{definition}
\newtheorem{defn}{Définition}

\theoremstyle{definition}
\newtheorem*{nota}{Notation}

\theoremstyle{definition}
\newtheorem*{limi}{Liminaire}

\theoremstyle{remark}
\newtheorem{rmq}{Remarque}

\theoremstyle{plain}
\newtheorem*{fait}{Fait}


%Ensembles de nombres
\newcommand{\N} {\mathbb{N}}
\newcommand{\Ne}{\N^\ast}
\newcommand{\Z} {\mathbb{Z}}
\newcommand{\Q} {\mathbb{Q}}
\newcommand{\R} {\mathbb{R}}
\newcommand{\Rb}{\overline{\mathbb{R}}}
\newcommand{\Rp}{\R_+}
\newcommand{\Rm}{\R_-}
\newcommand{\K} {\mathbb{K}}
\newcommand{\Ket} {\mathbb{K^{\ast}}}
\newcommand{\Cx}{\mathbb{C}}

%Ensembles, tribus
\newcommand{\A}{\mathbb{A}}
\renewcommand{\P}{\mathcal{P}}
\newcommand{\T}{\mathcal{T}}
\newcommand{\F}{\mathcal{F}}
\newcommand{\C} {\mathcal{C}}
\newcommand{\ouv}{\mathcal{O}}
\newcommand{\B}{\mathcal{B}}
\newcommand{\M}{\mathcal{M}}
\newcommand{\etag}{\mathcal{E}}
\newcommand{\etagOp}{\etag^{+}(\Omega)}
\newcommand{\etagO}{\etag(\Omega)}
%%Lp avant quotientage
\newcommand{\Lic}{\mathcal{L}^1}
\newcommand{\Lpc}{\mathcal{L}^p}
\newcommand{\Linfc}{\mathcal{L}^{\infty}}
\newcommand{\Lifull}{\Lic(\Omega, \B, m)}
\newcommand{\Lpfull}{\Lpc(\Omega, \B, m)}
\newcommand{\Linffull}{\Linfc(\Omega, \B, m)}
%%Lp après quotientage
\newcommand{\Li}{L^1}
\newcommand{\Ld}{L^2}
\newcommand{\Lp}{L^p}
\newcommand{\Linf}{L^{\infty}}
\newcommand{\Listd}{\Li(\Omega, m)} 
\newcommand{\Ldstd}{\Ld(\Omega, m)} 
\newcommand{\Lpstd}{\Lp(\Omega, m)} 

%Quelques opérateurs
\DeclareMathOperator{\codim}{codim}
\DeclareMathOperator{\rg}{rg}
\DeclareMathOperator{\tr}{tr}
\DeclareMathOperator{\vect}{Vect}
\DeclareMathOperator{\ima}{Im}
\DeclareMathOperator{\id}{Id}
\DeclareMathOperator{\sign}{sign}
\DeclareMathOperator{\ext}{ext}
\newcommand{\ind}{\mathbbm{1}}
\newcommand*{\spr}[2]{\left(\,#1\,\middle|\,#2\,\right)}
\newcommand{\interior}[1]{%
  {\kern0pt#1}^{\mathrm{o}}%
}
\DeclareMathOperator{\card}{card}
\DeclareMathOperator{\ppcm}{ppcm}

%Commandes de pure flemme
\newcommand{\veps}{\varepsilon}
\newcommand{\intab}{\int_{a}^{b}}
\newcommand{\intR}{\int_{-\infty}^{\infty}}
\newcommand{\intinf}{\int_{- \infty}^{+ \infty}}
\newcommand{\equi}{\Leftrightarrow}
\newcommand{\Kev}{$\K$-espace vectoriel }
\newcommand{\Kevs}{$\K$-espaces vectoriels }
\newcommand{\Rev}{$\R$-espace vectoriel }
\newcommand{\Revs}{$\R$-espaces vectoriels }
\newcommand{\Cev}{$\Cx$-espace vectoriel }
\newcommand{\Cevs}{$\Cx$-espaces vectoriels }
\newcommand{\evn}{(E, \norm{.})}
\newcommand{\interf}[2]{[#1,#2]}
\newcommand{\limb}{\lim\limits_}
\newcommand{\lebn}{\lambda_n}
\newcommand{\pp}{\text{~presque partout}}
\newcommand{\derivp}[2]{\frac{\partial #1}{\partial #2}}
\newcommand{\famiu}{(u_i)_{i \in I}}
\newcommand{\sev}{\text{~s.e.v.}}
\newcommand{\Mnp}{M_{n,p}(\K)}
\newcommand{\Mn}{M_{n}(\K)}
\newcommand{\tq}{\text{~t.q.~}}
\newcommand{\mq}{~montrer~que~}
\newcommand{\et}{\quad \text{et} \quad}
\newcommand{\ou}{\text{~ou~}}
\newcommand{\Kx}{\K[X]}


%Commandes de gars chelou
\renewcommand{\subset}{\subseteq}

%Commande restriction, ty duckduckgo
\def\restr#1#2{\mathchoice
              {\setbox1\hbox{${\displaystyle #1}_{\scriptstyle #2}$}
              \restraux{#1}{#2}}
              {\setbox1\hbox{${\textstyle #1}_{\scriptstyle #2}$}
              \restraux{#1}{#2}}
              {\setbox1\hbox{${\scriptstyle #1}_{\scriptscriptstyle #2}$}
              \restraux{#1}{#2}}
              {\setbox1\hbox{${\scriptscriptstyle #1}_{\scriptscriptstyle #2}$}
              \restraux{#1}{#2}}}
\def\restraux#1#2{{#1\,\smash{\vrule height .8\ht1 depth .85\dp1}}_{\,#2}} 

%Quelques normes
\DeclarePairedDelimiter\abs{\lvert}{\rvert}
\DeclarePairedDelimiter\labs{\bigl\lvert}{\bigr\rvert}
\DeclarePairedDelimiter\norm{\lVert}{\rVert}
\DeclarePairedDelimiterXPP\normi[1]{}\lVert\rVert{_1}{\ifblank{#1}{\:\cdot\:}{#1}}
\DeclarePairedDelimiterXPP\normd[1]{}\lVert\rVert{_2}{\ifblank{#1}{\:\cdot\:}{#1}}
\DeclarePairedDelimiterXPP\normp[1]{}\lVert\rVert{_p}{\ifblank{#1}{\:\cdot\:}{#1}}
\DeclarePairedDelimiterXPP\normq[1]{}\lVert\rVert{_q}{\ifblank{#1}{\:\cdot\:}{#1}}
\DeclarePairedDelimiterXPP\norminf[1]{}\lVert\rVert{_\infty}{\ifblank{#1}{\:\cdot\:}{#1}}

%Bracket en angle
\DeclarePairedDelimiter\angl{\langle}{\rangle}

%Set, parenthèses
\newcommand{\set}[1]{\{ #1 \}}
\newcommand{\bigset}[1]{\bigl\{ #1 \bigr\}}
\newcommand{\Bigset}[1]{\Bigl\{ #1 \Bigr\}}
\newcommand{\bigpar}[1]{\bigl( #1 \bigr)}
\newcommand{\Bigpar}[1]{\Bigl( #1 \Bigr)}

%Opitmisation des opérateurs de comparaison
\DeclarePairedDelimiter\tio{o (}{)}
\DeclarePairedDelimiter\gro{O (}{)}

\renewcommand{\F}{\mathbb{F}}
\newcommand{\U}{\mathbb{U}}

\pagestyle{fancy}

\fancyhead[C]{}
\fancyhead[R]{}

\begin{document}

\underline{Source :} Rombaldi - Algèbre - page 382.

\underline{Leçons :}
\begin{itemize}
\item 
\end{itemize}

\begin{nota}
Pour $p \geq 2$, on note $\F_p = \Z/p\Z$.
On note aussi $\U_n = \set{z \in \Cx ; z^n = 1}$ l'ensemble des racines $n$-ièmes de l'unité.
$\U_n$ est aussi le groupe cyclique engendré par $\omega_1 = \exp \left ( \dfrac{2i \pi}{n} \right )$.
Pour $\omega \in \U_n$, on désigne par $P_{\omega}$ le polynôme minimal sur $\Q$ de $\omega$,
i.e. le polynôme unitaire de $\Q[X]$ de degré minimal qui annule $\omega$. Ce polynôme est irréductible.
On appelle racine primitive $n$-ième de l'unité tout générateur de $\U_n$ et on note $R_n$ l'ensemble
de ces racines primitives.
\end{nota}

\section{Le critère d'Eisenstein}

\begin{defn}[Contenu]
Soit $P$ un polynôme non nul de $\Z[X]$ on appelle contenu de $P$, et on note $c(P)$ le pgcd des coefficients de $P$.
Si le contenu de $P$ est 1, on dit que $P$ est primitif.
\end{defn}

\begin{lem}
\label{lem:1}
Le produit de deux polynômes primitifs non nuls est primitif.
\end{lem}

\begin{proof}
Soient $P, Q$ tels qu'énoncés et $R$ leur produit. Supposons que $R$ n'est pas primitif, il existe donc un nombre premier
$p$ qui divise tous les coefficients de $R$. Autrement dit $\overline{R} = \overline{0}$ dans l'anneau intègre $\F_p[X]$
ce qui implique que $\overline{P} = \overline{0}$ ou $\overline{Q} = \overline{0}$, ce qui est absurde.
\end{proof}

\begin{lem}
\label{lem:2}
Pour $P, Q \in \Z[X] \setminus \set{0}$, on a $c(PQ) = c(P)c(Q)$.
\end{lem}

\begin{proof}
On écrit :
\[ PQ = c(P) \left ( \frac{P}{c(P)} \right ) c(Q) \left ( \frac{Q}{c(Q)} \right ) , \]
avec $\left ( \frac{P}{c(P)} \right )$ et $\left ( \frac{Q}{c(Q)} \right )$ primitifs, autrement dit :
\[ PQ = c(P) c(Q) R , \]
avec $R$ primitif selon le lemme précédent. On conclut par homogénéité du pgcd.
\end{proof}

\begin{lem}
\label{lem:3}
Soient $P, Q \in \Q[X]$ deux polynômes unitaires tels que $PQ \in \Z[X]$ alors $P$ et $Q$ sont aussi dans $\Z[X]$.
\end{lem}

\begin{proof}
Notons $m$ un dénominateur commun à tous les coefficients de $P$ et $Q$, on a alors $P_1=mP$ et $Q_1=mQ$ deux polynômes 
de $\Z[X]$ de coefficients dominants égaux à $m$.
On a alors $m^2 PQ = P_1 Q_1$ et donc $c(m^2 PQ)= m^2 c(PQ) = m^2$ car $P$ et $Q$ sont unitaires et $\Z$ intègre.
On déduit que $c(P_1)c(Q_1) = m^2$ selon le lemme précédent avec $c(P_1)$ et $c(Q_1)$ qui divisent $m$.
On a donc $c(P_1)=c(Q_1)=m$, i.e. $m$ divise tous les coefficients de $P$ et $Q$, 
donc $P=\dfrac1{m} P_1$ et $Q=\dfrac1{m} Q_1$ sont dans $\Z[X]$.
\end{proof}

\begin{thm}
Soit $P(X) = \sum_{k=0}^n a_k X^k$ un polynôme non constant dans $\Z[X]$ de degré $n \geq 2$. 
S'il existe un nombre premier $p$ tel que :
\begin{enumerate}[label=(\roman*)]
\item $p$ divise les coefficients $a_0, a_1, \dots, a_{n-1}$ ;
\item $p$ ne divise par $a_n$ ;
\item $p^2$ ne divise pas $a_0$ ;
\end{enumerate}
alors le poynôme $P$ est irréductible dans $\Q[X]$.
\end{thm}

\begin{proof}
Suppons le polynôme $P$ tel qu'énoncé mais réductible dans $\Q[X]$. Alors il est réductible dans $\Z[X]$ selon le lemme précédent.
On a donc $P=QR$ avec $Q, R$ des polynômes non constants dans $\Z[X]$ et on note :
\[ Q(X) = \sum_{k=0}^q b_k X^k \et R(X) = \sum_{k=0}^r c_k X^k , \]
avec les $b_k$ et les $c_k$ dans $\Z$ et $1 \leq q, r \leq n-1$.
Selon les hypothèses (i) et (ii) on a :
\[ \overline{P} = \overline{Q} \overline{R} = \overline {a_n} X^n \neq 0 , \]
dans $\F_p[X]$. Par ailleurs, $\overline{a_n} = \overline{b_q} \overline{c_r}$ dans le corps $\F_p$, donc ni $b_q$ ni $c_r$ ne sont nuls.
Comme la décomposition en facteurs irréductibles est unique (à l'ordre près) dans $\F_p[X]$, 
l'égalité $\overline{P} = \overline{b_q} \overline{c_r} X^n$ implique que :
\[ \overline{Q}(X) = \overline{b_q} X^q \et \overline{R}(X) = \overline{c_r} X^r . \]
Autrement dit, $p$ divise tous les coefficients (sauf dominants) de $Q$ et $R$, 
en particulier $p^2$ divise $b_0 c_0 = a_0$ ce qui est contradictoire avec l'hypothèse (iii).
\end{proof}

\section{Application aux polynômes cyclotomiques}

\begin{lem}
\label{lem:4}
La racine $n$-ième de l'unité $\omega_k = \exp \left ( \dfrac{2ik \pi}{n} \right )$ engendre $\U_n$ si, et seulement si, $k$ et $n$
sont premiers entre eux.
\end{lem}

\begin{proof}
Si $\omega_k$ engendre $\U_n$, il existe $p \in \N$ tel que $\omega_k^p = \omega_1 ^{kp} = \omega_1$.
Alors $\omega_1^{1-kp} = 1$ et on déduit que $n$ divise $1-kp$, il existe donc $q$ tel que $1= nq + kp$ donc $n$ et $p$ sont
premiers entre eux. \\
Réciproquement, si $n$ et $p$ sont premiers entre eux, on écrit $nq + kp = 1$ et on a 
\[ \omega_k^{p} = \omega_1^{1-nq} = \omega_1 , \]
et donc $\U_n = \angl{ \omega_k}$.
\end{proof}

\begin{coro}
On a $\card(R_n) = \varphi(n)$ avec $\varphi$ la fonction indicatrice d'Euler.
\end{coro}

\begin{lem}
\label{lem:5}
Soit $p$ un nombre premier.
\begin{enumerate}
\item Si $A_1, \dots, A_k$ sont des polynômes dans $\Z[X]$ et $A = \sum_{i=1}^k A_i$, on a alors l'égalité suivante dans $F_p[X]$ :
\[ \overline{A^p} = \sum_{i=1}^k \overline{A_i^k} . \]
\item Si $A$ est un polynôme dans $\Z[X]$, on a alors  l'égalité suivante dans $F_p[X]$ :
\[ \overline{A^p}(X) = \overline{A}(X^p) . \]
\end{enumerate}
\end{lem}

\begin{proof}
\begin{enumerate}
\item La résultat est vrai pour $k=1$. Pour $k=2$, considérons $A^p=(A_1+A_2)^p$.
Comme $p$ divise tous les coefficients binomiaux $\binom{p}{i}$ pour $1 \leq i \leq p-1$ on a :
\[ \overline{A^p} = \overline{A_1^p} + \overline{A_2^p}, \]
dans $\F_p$. On peut donc conclure par récurrence. 
\item On écrit $A(X) = \sum a_k X^k$, en utilisant le résultat précédent, on a :
\[ \overline{A^p}(X) = \sum \overline{a_k^p} X^{kp} . \]
D'après le théorème de Fermat, $\overline{a_k^p} = \overline{a_k}^p = \overline{a_k}$, et donc :
\[ \overline{A^p}(X) = \sum \overline{a_k} (X^{p})^k = \overline{A}(X^p) . \]
\end{enumerate}
\end{proof}

\begin{lem}
\label{lem:6}
Pour tout $\omega \in \U_n$, le polynôme minimal $P_\omega$ est dans $\Z[X]$.
\end{lem}

\begin{proof}
Soit $A \in \Q[X]$ s'annulant en $\omega$. Par division euclidienne, il est clair que $A$ est multiple de $P_\omega$ dans $\Q[X]$.
En particulier, $X^n-1$ est multiplie de $P_\omega$ soit $X^n-1 = P_\omega(X) Q(X)$ avec $Q \in Q[X]$.
$Q$ est nécessairement unitaire, donc d'après le lemme $\ref{lem:3}$ on déduit que  $P_\omega$ (et $Q$) sont dans $\Z[X]$.
\end{proof}

\begin{defn}
Pour $n \in \Ne$, on appelle polynôme cyclotomique le polynôme :
\[ \phi_n(X) = \prod_{z \in R_n} X-z . \]
\end{defn}

\begin{thm}
Pour tout $n \in \Ne, \phi_n$ est irréductible dans $\Q[X]$.
\end{thm}

\begin{proof}
Soit $\omega \in R_n, p$ un nombre premier ne divisant pas $n$ et on note $Q_{\omega}$ le polynôme minimal sur $\Q$
de $\omega^p$. Si $Q_{\omega}$ et $P_{\omega}$ sont différents, comme ils sont irréductibles dans $\Q[X]$, il sont donc premiers
entre eux dans $\Q[X]$. 
Il en résulte que le polynôme $X^n-1$ qui s'annule en $\omega$ et en $\omega_p$ est multiple de $Q_{\omega}$ et $P_{\omega}$ 
donc du produit $Q_{\omega} P_{\omega}$.
On a donc ,
\[ X^n-1 = Q_{\omega}(X) P_{\omega}(X) Q(X) , \]
avec $Q(x) \in \Z[X]$ selon le lemme $\ref{lem:3}$  
(le polynôme $Q_{\omega} P_{\omega}$ est unitaire dans $\Q[X]$ et $Q$ est nécessairement unitaire). 
Le polynôme $Q_{\omega}(X^p)$ s'annule en $\omega$, il est donc multiple de $P_{\omega}$ et s'écrit :
\[ Q_{\omega}(X^p) = P_{\omega}(X) R(X) , \]
avec $R \in \Q[X]$ nécessairement unitaire et $Q_{\omega} \in \Z[X]$ (cf. lemme $\ref{lem:6}$) 
donc $R \in \Z[X]$ (cf. lemme $\ref{lem:3}$).
On se place maintenant dans $\F_p[X]$ et on a, selon le lemme $\ref{lem:5}$ :
\[ \overline{Q_\omega^p}(X) = \overline{Q_{\omega}}(X^p) = \overline{P_{\omega}}(X) \overline{R}(X) . \]
Prenons $\overline{S}$ un polynôme irréductible de $\F_p[X]$ qui divise $\overline{\P_\omega}$. 
Alors $\overline{S}$ divise $\overline{Q_\omega^p} = \overline{Q_\omega}^p$, et donc, par irréductibilité, 
il divise $\overline{Q_\omega}$. En conséquence, $\overline{S}^2$ divise $X^n - \overline{1}$
(car ce dernier polynôme est produit de $\overline{\P_\omega}$ et $\overline{\Q_\omega}$).
On peut donc écrire :
\[  X^n - \overline{1} = \overline{S}(X) \overline{T}(X) , \]
avec $T$ un polynôme. En dérivant formellement, on obtient :
\[ \overline{n} X^{n-1} = \overline{S}(X) \overline{U}(X) , \]
avec $U$ un polynôme. 
On déduit que si $\overline{S}$ est irréductible dans $\F_p[X]$ alors il est unitaire et divise $\overline{n} X^{n-1}$.
De plus, comme $p$ ne divise pas $n$, on a $\overline{n} \neq 0$, donc on a nécessairement $\overline{S} = X$ ou $1$.
Si $\overline{S} = X$ alors $\overline{S}$ divise $X - \overline{1}$ est contradictoire et donc $\overline{S} = 1$.
Ceci implique que $\overline{P_\omega}$ est de degré nul ce qui est impossible.
Donc, l'hypothèse de départ, à savoir $P_\omega \neq Q_\omega$ est fausse, donc $P_\omega = Q_\omega$. \\

Prenons $\omega = \omega_1$, on va conclure en montrant que $P_\omega$ s'annule pour tout $\omega_k \in R_n$. 
Si $\omega_k \in R_n \subset \U_n$, alors $k$ est premier avec $n$ (et $k \in \set{1, \dots, n}$).
En écrivant la décomposition en facteurs premiers de $k, k =\prod p_i$ (avec les $p_i$ non nécessairement distincts),
on sait que chaque $p_i$ est premier avec $n$. On a donc :
\[ P_{\omega} = P_{\omega^{p_1}} , \]
et ainsi, de proche en proche :
\[ P_{\omega} = P_{\omega^{p_1 p_2}} = \dots = P_{\omega^{k}} , \]
et donc $P_{\omega}(\omega^k) =  P_{\omega^{k}}(\omega^k) = 0$ pour tout $\omega^k \in R_n$.
On a donc $P_{\omega} = \phi_n$ et ce premier étant irréductible, on déduit le résultat.
\end{proof}

\end{document}